% =====================================================================
\section{The whole machine on one page}\label{sec:architecture}
% =====================================================================

Figure~\ref{fig:pipeline} is the complete discovery pipeline as finally assembled in Phase~6. Read it top to
bottom: each stage hands its output to the next. The right-hand column says which package does the work and
which question the stage answers.

\begin{figure}[!htb]
\centering
% flowchart 'pipeline' -- used by docs/Report.tex
\begin{tikzpicture}[node distance=3.2mm]
  \tikzset{st/.style={draw=#1!75!black, fill=#1!9, rounded corners=3pt, line width=0.7pt,
                      font=\small, align=left, inner sep=4pt, text width=7.6cm},
           side/.style={font=\scriptsize, align=left, text width=5.6cm, text=black!75}}
  \node[st=cData] (s1) {\textbf{1. Observations}\quad 400 noisy points $(x,t,u)$; no equation given};
  \node[st=cRag, below=of s1] (s2) {\textbf{2. Describe + retrieve (RAG)}\quad turn data statistics into
        text; fetch relevant textbook chunks};
  \node[st=cKg, below=of s2] (s3) {\textbf{3. Knowledge graph}\quad documented mechanisms $\to$ equations,
        parameters, sources};
  \node[st=cLlm, below=of s3] (s4) {\textbf{4. Generate hypotheses (LLM)}\quad schema-validated JSON
        equations, grounded in the evidence};
  \node[st=cHyp, below=of s4] (s5) {\textbf{5. Candidate pool + filters}\quad LLM ideas $\cup$ KG templates;
        deduplicate; syntax, variables, physics pre-checks};
  \node[st=cPinn, below=of s5] (s6) {\textbf{6. Inverse PINN per candidate}\quad compile the equation, train
        on the train split, estimate parameters};
  \node[st=cVal, below=of s6] (s7) {\textbf{7. Gate on held-out data}\quad reject unless held-out error,
        physics, BC, IC and bounds all pass};
  \node[st=cVal, below=of s7] (s8) {\textbf{8. Model selection}\quad term necessity $\to$ fit-equivalence
        $\to$ lowest complexity $\to$ BIC};
  \node[st=cVal, below=of s8] (s9) {\textbf{9. OOD validation}\quad predict $t>0.8$ observations never used
        for any decision};
  \node[st=cVal, below=of s9] (s10) {\textbf{10. Novel prediction}\quad solve the selected law forward beyond
        the data ($t$ up to 1.5)};
  \foreach \a/\b in {s1/s2,s2/s3,s3/s4,s4/s5,s5/s6,s6/s7,s7/s8,s8/s9,s9/s10}
     \draw[flow] (\a) -- (\b);
  \node[side, right=4mm of s1]  {\fn{data/}\\ ``What did we measure?''};
  \node[side, right=4mm of s2]  {\fn{evaluation/describe.py}, \fn{rag/}\\ ``What is documented about such behaviour?''};
  \node[side, right=4mm of s3]  {\fn{knowledge_graph/}\\ ``How do documented concepts relate?''};
  \node[side, right=4mm of s4]  {\fn{llm/}, \fn{hypotheses/pipeline.py}\\ ``Which equations are worth testing?''};
  \node[side, right=4mm of s5]  {\fn{hypotheses/search.py}, \fn{validation.py}, \fn{constraints.py}\\ ``Is each one well-formed and physical?''};
  \node[side, right=4mm of s6]  {\fn{hypotheses/adapter.py}, \fn{pinn/}\\ ``Can data + this physics be fitted together?''};
  \node[side, right=4mm of s7]  {\fn{hypotheses/search.py}\\ ``Does it explain unseen observations?''};
  \node[side, right=4mm of s8]  {\fn{validation/model_selection.py}\\ ``Which is the simplest sufficient law?''};
  \node[side, right=4mm of s9]  {\fn{hypotheses/search.py}\\ ``Does it extrapolate in time?''};
  \node[side, right=4mm of s10] {\fn{validation/forward_solver.py}\\ ``What happens where no data exists?''};
\end{tikzpicture}

\caption{The end-to-end discovery pipeline (Phase~6 form). Colours follow the package legend used in every
flowchart.}
\label{fig:pipeline}
\end{figure}

\begin{lesson}[: who is allowed to decide]
\begin{itemize}
  \item RAG \emph{retrieves}, the graph \emph{organises}, the LLM \emph{proposes}. None of them can mark an
        equation as correct; there is literally no field for that in the schema.
  \item Ranking (\fn{hypotheses/ranking.py}) only sets the \emph{order} in which candidates are examined.
  \item Only measured numbers decide: the gate (\fn{validate_candidate}) and \fn{select_model}, both using
        held-out \emph{noisy} observations, never the hidden truth.
  \item Status strings are deliberately modest: ``supported by the available observations and physics
        constraints'' or ``rejected under the tested conditions''. Never ``proven''.
\end{itemize}
\end{lesson}

\subsection{How the code packages depend on each other}

Figure~\ref{fig:packages} was computed from the actual \fn{import} statements in the code. An arrow
$A\to B$ means ``package $A$ imports from package $B$''. Double-headed arrows are mutual imports.

\begin{figure}[!htb]
\centering
% flowchart 'packages' -- used by docs/Report.tex
\begin{tikzpicture}[
   pk/.style={draw=#1!75!black, fill=#1!12, rounded corners=4pt, line width=0.8pt,
              font=\ttfamily\small\bfseries, minimum width=3.3cm, minimum height=9mm, align=center},
   dep/.style={->, line width=0.7pt, black!60},
   mdep/.style={<->, line width=0.7pt, black!60}]
  \node[pk=cRag]  (rag) at (0,0)     {rag\\[-1pt]{\normalfont\scriptsize Phase 5}};
  \node[pk=cKg]   (kg)  at (0,-2.4)  {knowledge\_graph\\[-1pt]{\normalfont\scriptsize Phase 5}};
  \node[pk=cLlm]  (llm) at (5,0)     {llm\\[-1pt]{\normalfont\scriptsize Phase 4}};
  \node[pk=cHyp]  (hyp) at (5,-2.4)  {hypotheses\\[-1pt]{\normalfont\scriptsize Phases 4, 6}};
  \node[pk=cVal]  (val) at (5,-4.8)  {validation\\[-1pt]{\normalfont\scriptsize Phases 3, 6, 7}};
  \node[pk=cPinn] (pinn) at (0,-4.8) {pinn\\[-1pt]{\normalfont\scriptsize Phase 3}};
  \node[pk=cData] (data) at (0,-7.2) {data\\[-1pt]{\normalfont\scriptsize Phase 1}};
  \node[pk=cEqd]  (eqd) at (10,-4.8) {equation\_discovery\\[-1pt]{\normalfont\scriptsize Phase 2}};
  \node[pk=cEval] (ev)  at (10,-1.2) {evaluation\\[-1pt]{\normalfont\scriptsize Phase 7}};
  \node[pk=cExp]  (exp) at (5,-7.2)  {experiments\\[-1pt]{\normalfont\scriptsize 02--11 scripts}};

  \draw[mdep] (rag) -- (kg);
  \draw[mdep] (pinn) -- (data);
  \draw[mdep] (hyp) -- (val);
  \draw[dep] (hyp) -- (kg);
  \draw[dep] (hyp) -- (llm);
  \draw[dep] (hyp) -- (pinn);
  \draw[dep] (val) -- (pinn);
  \draw[dep] (val) -- (eqd);
  \draw[dep] (val.east) .. controls +(1.6,1.2) and +(1.0,-0.6) .. (llm.south east);
  \draw[dep, cEval!80!black] (ev) -- (llm);
  \draw[dep, cEval!80!black] (ev) -- (hyp);
  \draw[dep, cEval!80!black] (ev) -- (eqd);
  \draw[dep, cEval!80!black] (ev) -- (val);
  \draw[dep, cEval!80!black] (ev.north) .. controls +(0,1.6) and +(1.5,1.2) .. (rag.north east);
  \draw[dep, cEval!80!black, rounded corners=6pt] (ev.west) -- (2.4,-1.2) -- (kg.north east);
  \node[font=\scriptsize, align=center, text width=4.4cm, below=2mm of exp]
       {\fn{experiments/} imports all nine packages; \fn{tests/} imports what each phase tests};
\end{tikzpicture}

\caption{Package dependency map, derived from the real imports. The phase that introduced each package is shown
under its name.}
\label{fig:packages}
\end{figure}

\begin{intuition}[: reading the map]
\fn{hypotheses} sits in the middle because it glues the brainstormer (\fn{llm}) to the lab (\fn{pinn}) and the
judge (\fn{validation}). \fn{evaluation} (Phase~7) reaches into almost everything because it re-runs the whole
pipeline on benchmark datasets. \fn{data} and \fn{pinn} import each other only because the problem description
and the PINN share the same data-loading helper (\fn{load_observational_data}).
\end{intuition}
